\chapter{Discussion}

\section{Summary of Findings}

The empirical results support three main observations. First, the image-only DenseNet121 model achieves the highest discrimination and threshold-based performance on the held-out temporal test set, while the clinical models are substantially weaker on AUROC and AUPRC. This indicates that, within the present cohort and label definition, chest radiographs contain substantially more task-relevant information for radiographic pneumonia prediction than the selected triage variables alone.

Second, early multimodal fusion does not improve over the image-only baseline. Although the multimodal model remains competitive in absolute terms and clearly exceeds the clinical-only baselines, point estimates favor the image model, and paired bootstrap intervals for multimodal-minus-image differences include zero for both AUROC and AUPRC. The fraction of patient-level bootstrap resamples favoring the multimodal model is also low. Taken together, these results do not support the claim that multimodal fusion provides a reliable performance advantage in the present setup.

Third, the auxiliary analyses show that model behavior is more nuanced than a single discrimination metric suggests. Calibration differs from discrimination, with the multimodal model exhibiting somewhat better-calibrated probabilities than the image-only model under fixed-bin calibration metrics despite weaker ranking performance. Decision curve analysis places both deep models above the clinical baselines, but again does not establish a clear multimodal advantage across thresholds. Finally, performance varies substantially across negative-case subtypes, with marked degradation when the negative class is restricted to other abnormal radiographs rather than relatively normal studies. This indicates that task difficulty depends strongly on the structure of the comparator group.

Taken together, these findings support a central conclusion of the thesis: under strict temporal constraints and a radiograph-centered label definition, a strong image model dominates predictive performance, while multimodal fusion with triage variables does not provide a consistent additional benefit.

\section{Interpretation of Multimodal Performance}

The absence of a multimodal advantage is one of the most important results of this thesis. While multimodal learning is often assumed to be beneficial, the present findings suggest that such benefit is conditional rather than automatic. Several plausible factors may explain why the multimodal model does not improve on the image-only baseline, although the current experiments cannot establish any of them as definitive causal mechanisms.

A first factor is the limited incremental value of the available tabular data. The structured clinical branch is restricted to triage variables available at or before the imaging time. These measurements encode patient state at presentation, but they are relatively coarse compared with the spatially rich information present in a chest radiograph. Once an image has already been observed by the model, triage variables may add only limited independent signal for the specific task of radiographic pneumonia prediction.

A second factor is redundancy between modalities. Acute respiratory illness can manifest simultaneously in both physiology and imaging. Tachypnea, low oxygen saturation, and elevated temperature may correlate with pneumonia, but they may also correlate with many other respiratory or systemic conditions. If these variables overlap strongly with information that is already implicitly represented in the radiograph embedding, then early fusion may add complexity without adding enough unique signal to improve discrimination.

A third factor is representation-level asymmetry between the two branches. In the implemented multimodal model, the image encoder produces a high-dimensional visual representation learned from multilabel pretraining and task-specific fine-tuning, whereas the structured branch is built from a relatively small set of triage variables. It is therefore plausible that the fused representation is dominated by the image pathway, with the clinical branch contributing only weakly or inconsistently. This does not mean the triage information is irrelevant in principle, but it does suggest that its incremental value may be limited under the present task formulation.

A fourth factor is architectural limitation. The implemented multimodal model uses a relatively simple early-fusion strategy: a DenseNet-derived image embedding is concatenated with a learned tabular representation and passed through a compact fusion head. This is a reasonable and interpretable design for a thesis, but it may not fully exploit complex cross-modal relationships. More expressive fusion mechanisms, alternative conditioning schemes, or deeper tabular encoders might behave differently. However, the thesis result remains meaningful even without such extensions: under the chosen and well-defined fusion design, no robust multimodal gain is observed.

A fifth factor is optimization difficulty. Multimodal systems are harder to train than unimodal ones because they require the coordinated optimization of multiple branches and a fusion mechanism. Even if the tabular modality contains some complementary signal, that signal may be difficult to exploit effectively without more careful architectural or optimization choices. In that sense, a negative result for early fusion does not imply that clinical information is irrelevant in principle; rather, it shows that its benefit is not guaranteed in the present formulation.

For a thesis, this finding is valuable precisely because it is honest and constrained. The result is not framed as a failure of multimodal learning in general, but as evidence that multimodal benefit must be demonstrated empirically within a carefully defined setting.

\section{Role of Temporal Constraints}

A defining methodological feature of this work is the explicit anchoring of prediction at the imaging time, with strict exclusion of features that may depend on post-imaging events. This decision is central to the scientific validity of the thesis. In medical machine learning, strong performance can be misleading when the feature set contains variables that are not genuinely available at prediction time. Such leakage may arise subtly, especially in multimodal pipelines that combine data sources recorded at different stages of the clinical workflow.

In the present study, triage inputs are aligned with a pre-result clinical state, and variables such as disposition are excluded due to their potential dependence on downstream decisions. Likewise, patient-level temporal splitting prevents the same patient from appearing across train and test in a way that would compromise the interpretation of generalization. Together, these decisions make the reported performance more conservative, but also more credible.

This point is important for interpreting the absolute metric values. A less restrictive pipeline could potentially yield higher AUROC or AUPRC by exploiting future information, repeated-patient structure, or downstream documentation artifacts. However, such gains would not reflect a clean solution to the intended prediction problem. The current setup instead prioritizes temporal validity over maximal feature inclusion, ensuring that reported performance corresponds to a well-defined clinical decision point.

From a thesis perspective, this design choice is a strength. Even if temporal rigor makes the problem harder and reduces the apparent payoff of multimodal fusion, it increases the trustworthiness of the conclusions. The thesis therefore emphasizes methodological validity rather than maximizing benchmark performance at all costs.

\section{Error Structure and Negative Definition}

The abnormal-negative analysis provides one of the clearest insights into the structure of the task. Stratification by CheXpert abnormality labels reveals a substantial decline in performance when negatives are restricted to studies with other radiographic abnormalities, compared with relatively normal negatives. This suggests that the main challenge is not simply to distinguish pneumonia from normal radiographs, but to distinguish pneumonia from other abnormal thoracic presentations that may share overlapping visual patterns.

This result helps explain why the image-only model, despite being the strongest overall, still exhibits substantial error. A model may appear to perform well when the negative class is dominated by relatively normal cases, but its effective task becomes much harder when non-pneumonia abnormalities are included. In such settings, the decision boundary must separate among multiple visually similar pathological states rather than between clearly healthy and clearly diseased images.

The negative-definition result also supports the qualitative findings from Grad-CAM. In the present project, the interpretability examples are generated from the image-only DenseNet121 model and are intended as descriptive diagnostics rather than causal explanations. False positives appear to arise in cases where the model responds to visually abnormal but non-pneumonic structures or patterns, whereas false negatives arise when the radiographic signal is subtle, diffuse, or ambiguous. These observations are consistent with the idea that the learned image representation is highly sensitive to overt abnormality, but not perfectly specific for pneumonia as a distinct radiographic concept.

More broadly, this analysis illustrates that label design and evaluation cohort construction are inseparable from model interpretation. A binary pneumonia label derived from report structure is necessarily embedded within a broader ontology of thoracic findings. As a consequence, reported binary performance metrics must be interpreted with caution: they reflect not only model quality, but also the composition of positive and negative cases.

\section{Task Ambiguity as a Central Finding}

One of the most important conceptual findings of the thesis is that radiographic pneumonia detection, as implemented here, is an inherently ambiguous problem rather than a clean binary discrimination task. This ambiguity arises from at least three interacting sources.

The first source is label ambiguity. The primary target is derived from CheXpert report labels and therefore reflects a report-based radiographic proxy for pneumonia rather than a definitive etiologic diagnosis. The \texttt{u\_ignore} policy reduces ambiguity by excluding uncertain annotations, but it does not eliminate residual label noise or the broader mismatch between radiographic appearance and underlying disease mechanism. Even a carefully trained model is therefore learning against an imperfect target.

The second source is radiographic overlap. Pneumonia shares visual characteristics with other thoracic abnormalities, including pleural effusion, pulmonary edema, and non-specific parenchymal opacity. The hard-negative analysis makes this explicit: both image-only and multimodal performance decline substantially when the comparison class is restricted to other abnormal studies. This indicates that the difficulty of the task is not accidental or limited to a few edge cases, but is built into the visual structure of the problem itself.

The third source is limited specificity in the structured modality. Triage-at-\(t_0\) variables may reflect the presence of acute illness, but they are not highly specific for radiographic pneumonia. Fever, tachypnea, hypoxemia, and abnormal blood pressure can accompany many other conditions. As a result, adding these variables to the image pathway does not necessarily resolve the ambiguity created by overlapping radiographic findings.

Taken together, these three sources suggest that the key result of the thesis is not merely that image-only performance exceeds multimodal performance, but that the benchmark itself encodes an ambiguity-rich discrimination problem. Under this formulation, the image model is not failing because it is weak; rather, it is operating near the limits imposed by overlapping radiographic patterns, imperfect supervision, and limited complementary structured information. This framing helps explain both the hard-negative findings and the absence of a robust multimodal advantage.

\section{Calibration and Clinical Relevance}

Calibration metrics show that the deep models achieve lower Brier scores than the clinical baselines, indicating improved average agreement between predicted probabilities and observed outcomes. However, expected calibration error does not follow the same ranking as discrimination, with the image model exhibiting a higher point ECE than the multimodal model despite superior AUROC. Maximum calibration error also varies across models. These findings reinforce that discrimination and calibration measure different properties of predictive systems.

This distinction is especially relevant in clinical settings. A model with stronger ranking performance may still produce probabilities that are poorly aligned with empirical event frequencies, which could be problematic if outputs are interpreted as risk estimates rather than merely as ranking scores. The present results suggest that the multimodal model, despite weaker discrimination than the image-only system, may produce somewhat better-calibrated probabilities under the fixed-bin calibration metrics used here.

At the same time, lower ECE for the multimodal model does not by itself override the stronger AUROC and AUPRC of the image model. The preferred metric depends on the intended use. In ranking-oriented workflows, stronger discrimination may be more valuable; in settings where predicted probabilities are communicated directly or interpreted as risk estimates, calibration quality becomes more important. The thesis therefore does not claim a single universally superior model across all possible use cases, but rather shows that model preference depends on which property of predictive performance is prioritized.

Decision curve analysis offers a complementary view by asking whether a model would improve threshold-based decision making relative to simpler alternatives. In this study, both deep models achieve substantially higher net benefit than the clinical baselines across the evaluated thresholds. The relative ordering of the image-only and multimodal models is threshold-dependent, which is consistent with the broader finding that multimodal fusion is competitive but not clearly superior overall. This result is also consistent with the fixed-threshold main table: multimodal slightly improves precision at 0.5, whereas image retains better overall ranking and recall.

At the same time, these decision-oriented results should not be overinterpreted. Net benefit depends on assumptions about intervention trade-offs and threshold selection that are not estimated directly from this dataset. Thus, the DCA findings should be read as supportive evidence of relative utility rather than as a claim of immediate clinical readiness. More broadly, they reinforce the value of evaluating models beyond AUROC alone, especially when deployment decisions would depend on calibrated probabilities or threshold-based actions.

\section{Limitations}

Several limitations should be acknowledged. First, the study relies on CheXpert-derived pneumonia labels under a policy that excludes uncertain annotations. This reduces ambiguity relative to treating uncertain labels as negatives, but it may also remove challenging cases and does not eliminate residual label noise. The target remains a report-derived radiographic formulation of pneumonia rather than a definitive microbiological or clinically adjudicated diagnosis.

Second, all experiments are conducted on a single MIMIC-derived cohort. Although the dataset is large and widely used, it does not establish generalization to other hospitals, patient populations, imaging protocols, or clinical workflows. No external validation is performed in the present thesis, and therefore the conclusions should be understood as internally valid within the studied setting rather than universally generalizable.

Third, the multimodal analysis is intentionally restricted to triage features. This is methodologically consistent with the thesis emphasis on strict feature timing, but it also limits the breadth of the structured clinical modality. Additional sources such as laboratory data, medication history, or longitudinal context are not included in the primary comparison. Laboratory features were explored in the broader project, but their low coverage makes them more appropriate as a sensitivity analysis than as a core thesis modality.

Fourth, the architecture search space is limited. The multimodal model provides a clear and interpretable early-fusion baseline, but it is not an exhaustive search over possible multimodal designs. Similarly, the image model is strong but not necessarily optimal among all possible radiograph architectures. The goal of the thesis is rigorous comparative evaluation under controlled conditions, not maximal engineering optimization.

Fifth, most threshold-based summary metrics in the main results table are reported at a fixed probability threshold of 0.5. This is a standard reference point, but it should not be interpreted as an optimized clinical operating threshold. Different deployment scenarios could favor different thresholds and therefore different trade-offs between precision, recall, and net benefit.

Sixth, no post hoc recalibration comparison, such as temperature scaling or isotonic regression, is included in the primary analysis. As a result, the calibration results should be interpreted as properties of the raw model outputs rather than as the best achievable calibrated versions of these models.

Seventh, the abnormal-negative analysis is restricted to the subset of cases with available CheXpert abnormality annotations and should therefore be interpreted as a secondary descriptive analysis rather than as a separate benchmark. Even so, its consistency with the qualitative findings strengthens its interpretive value.

\section{Future Work}

Several directions for future work follow naturally from the present results. A first direction is to evaluate alternative target formulations, especially clinically actionable pneumonia definitions that differ from the current radiographic label. This would directly address the label ambiguity highlighted in the present study and could change the relative value of structured clinical inputs.

A second direction is to extend the multimodal modelling strategy itself. Richer fusion mechanisms, deeper tabular encoders, or alternative training schedules could be used to test whether the absence of multimodal gain is architectural rather than informational. In parallel, systematic ablations could be performed to identify which subsets of clinical variables contribute most, if at all, once radiographic information is available.

A third direction is to incorporate additional structured modalities under carefully defined timing rules. Laboratory data are an obvious candidate, but their usefulness depends on availability, timing, and missingness. In the current project, laboratory coverage is too limited for a fair primary comparison, but it remains a reasonable sensitivity-analysis direction and a possible extension in future datasets or expanded cohorts.

Finally, external and multi-site validation would be required to assess robustness and generalizability beyond the current MIMIC-based environment. Such validation is outside the scope of the thesis, but it represents an essential step for any future attempt to translate the conclusions into broader clinical use.